\scanpage{221}
Relations entre les 4 complexes, et fonctorialités.

Les \deleted{complexes} algèbres différentielles sont \unsure{fonctorielles|fonctoriels} en $(A,M)$. Pour $A$ fixé, $M$ variable, elles transforment sommes en produit tensoriel (d'où, pour chacune, application diagonale compatible avec \deleted{$d$ ou $\delta$} la différentielle), et commutent aux \uline{$\lim$} filtrantes.

L'hom can. d'algèbres (fonctoriel en $A,M$)
\diagram{\begin{tikzcd}[ampersand replacement=\&,column sep=large]
 \operatorname{Sym}^*M\otimes\Lambda^*M \arrow[r]\arrow[d,equal] \& \Gamma^*M\otimes\Lambda^*M \arrow[d,equal]\\
 C^*_{\mathrm{DR}}\operatorname{Sym}^*M \arrow[r] \& C^*_{\mathrm{DR\,p.d.}}(\Gamma^*M)
\end{tikzcd}}
est compatible avec les différentielles $d$. (C'est un isom \deleted{pour} si $A$ de car $0$ (réciproque vraie si \deleted{$A$} $M$ libre de rang $>0$), donc le calcul de coh. dans 1 est contenu dans celui de 3. [ce dernier se fait en réduisant au cas où \deleted{$A$}$M$ est libre de t.f. par Lazard, puis par Künneth au cas $M$ libre de rang $1$, auquel cas c'est trivial].

L'hom can. d'algèbres
\[
 \operatorname{Sym}^*M\otimes\Lambda^*M\unsure{\relbar|\longrightarrow}\Gamma^*M\otimes\Lambda^*M
\]
est aussi compatible avec les opérateurs diff. $\delta$. \deleted{\unsure{Comme} \unsure{lemme}} \deleted{\unsure{C'est}} Procédant comme dessus, le calcul d'homologie de (4) est ramené à celui de (2), qui peut se traiter par la $=$ méthode que pour (1).

La \uline{définition} de \unsure{l'op.|opér.} $\delta$ dans (2) et (4) peut d'ailleurs se ramener à un cas plus général de Koszul: si $S$ est un anneau comm, $M$ un $S$-module, \deleted{\unsure{libre}} et $\varphi\colon M\to S$ une forme

% archive sheet number: 220

